\section{The \texttt{Lens\_parab\_Cyl} McXtrace Component}
X-ray compound refractive lens (CRL) with a parabolic cylinder shape

\subsection*{Identification}
\begin{itemize}
  \item \textbf{Author:} Jana Baltser and Erik Knudsen
  \item \textbf{Origin:} NBI
  \item \textbf{Date:} April 2011
\end{itemize}

\subsection*{Description}
An X-ray compound refractive lens (CRL) with a parabolic cylinder profile focusing in 1D, i.e. onto a line

The lens is invariant along the x-axis and has a parabolic profile along the y-axis defined as z/c = y\textasciicircum{}2/b\textasciicircum{}2. Thus, i.e. it focuses onto a line along the x-axis. N\textgreater{}1 means that a stack of lenses is to be simulated. Each lens consists of a pair of two opposing parabolic surfaces with a distance d between them. The reference point of the component is at the bottom of the first parabolic surface.

Example: Lens\_parab\_Cyl(r=.5e-3,yheight=1.3e-3,xwidth=1.3e-3,d=.1e-3,N=21, material\_datafile="Be.txt")

\subsection*{Input parameters}
Parameters in \textbf{boldface} are required; the others are optional.

\begin{longtable}{p{0.22\textwidth}p{0.12\textwidth}p{0.46\textwidth}p{0.14\textwidth}}
\toprule
\textbf{Name} & \textbf{Unit} & \textbf{Description} & \textbf{Default} \\
\midrule
\endhead
material\_datafile & str & Datafile containing f1 constants & "Be.txt" \\
r & m & Radius of curvature (circular approximation at the tip of the profile). & .5e-3 \\
yheight & m & The CRL's aperture along Y. & 1.2e-3 \\
xwidth & m & The width of the CRL in the invariant direction x. & 1.2e-3 \\
d & m & Distance between two surfaces of the lens along the propagation axis. & .1e-3 \\
N & 1 & Number of single lenses in a stack. & 1 \\
rough\_z & rad & RMS value of random slope along z. & 0 \\
rough\_xy & rad & RMS value of random slope along x and y. & 0 \\
\bottomrule
\end{longtable}

\subsection*{Links}
\begin{itemize}
  \item Component source code found in file \texttt{Lens\_parab\_Cyl.comp}.
\end{itemize}
\IfFileExists{optics/Lens_parab_Cyl_static.tex}{\input{optics/Lens_parab_Cyl_static.tex}}{}